% ===================================================================
%  Семинар 2. Содержательная часть — единый источник для обоих файлов.
%
%  Состав задач взят из «Семинары.docx»:
%      основные:       2.8, 2.9, 2.16  — по «Семинар 2.docx», то есть
%                      то, что реально разбиралось на паре;
%      дополнительные: 2.23, 2.10, 2.11, 2.21, 2.13 — остаток списка
%                      семинара 2 из «Семинары.docx».
%      Источник: Сборник задач к начальному курсу эконометрики (РАНХиГС),
%                гл. 2 «Модель парной регрессии».
%
%  Порядок задач сохранён ровно тот, что указан в «Семинары.docx».
%  Задача 2.8 есть и в семинаре 1 (по «Семинары.docx» она отнесена туда),
%  и здесь: в семинаре 2 с неё начинается разбор регрессии без константы.
%
%  Раздел «Основные формулы» и всё внутри \begin{solution}...\end{solution}
%  попадает только в seminar_02_solutions.pdf. Блоки \begin{excel}...\end{excel}
%  печатаются в обоих файлах: это часть задания.
%
%  Условия задач — дословно по сборнику. Решения опираются на формулы
%  раздела 1 и ссылаются на них, не повторяя выводов.
% ===================================================================

\seminartitle{Семинар 2}{Суммы квадратов и $R^2$, свойства оценок,\\ теорема Гаусса--Маркова}

% ===================================================================
% Теория печатается только в файле с решениями.
\ifsolutions
\section{Основные формулы}

% -------------------------------------------------------------------
\subsection{Что используется из семинара 1}\label{sec:recap}

Модель и предпосылки классической регрессии те же, что в семинаре~1:
\begin{equation}
  Y_t = \alpha + \beta X_t + \varepsilon_t, \qquad t = 1,\dots,n,
  \label{eq:model}
\end{equation}
1)~выполнена спецификация \eqref{eq:model}; 2)~$X_t$ детерминированы и не все
равны между собой; 3a)~$\E\varepsilon_t = 0$ и $\D(\varepsilon_t) = \sigma^2$;
3b)~$\Cov(\varepsilon_t,\varepsilon_s) = 0$ при $t \neq s$;
3c)~$\varepsilon_t \sim N(0,\sigma^2)$. Обозначения прежние: $\ah$, $\bh$ —
МНК-оценки, $\yh_t = \ah + \bh X_t$ — подогнанные значения,
$e_t = Y_t - \yh_t$ — остатки, $x_t = X_t - \overline{X}$,
$y_t = Y_t - \overline{Y}$ — отклонения от средних.

Ниже используются как готовые следующие формулы семинара~1 (выводы — там же).
МНК-оценки и суммы в отклонениях:
\begin{gather}
  \bh = \frac{\sum x_ty_t}{\sum x_t^2},
  \qquad
  \ah = \overline{Y} - \bh\,\overline{X},
  \label{eq:ols} \\[2pt]
  \sum x_t^2 = \sum X_t^2 - \frac{\bigl(\sum X_t\bigr)^2}{n},
  \qquad
  \sum x_ty_t = \sum X_tY_t - \frac{\sum X_t\sum Y_t}{n}.
  \label{eq:dev-sums}
\end{gather}
Полная, объяснённая и остаточная суммы квадратов определены подписями в первой
строке \eqref{eq:ss}; эти определения годятся для любой регрессии. Само
равенство и формулы через отклонения во второй строке верны только для
регрессии со свободным членом:
\begin{equation}
\begin{gathered}
  \underbrace{\sum\bigl(Y_t - \overline{Y}\bigr)^2}_{\TSS}
  =
  \underbrace{\sum\bigl(\yh_t - \overline{Y}\bigr)^2}_{\RSS}
  +
  \underbrace{\sum\bigl(Y_t - \yh_t\bigr)^2}_{\ESS},
  \\[6pt]
  \TSS = \sum y_t^2,
  \qquad
  \RSS = \bh^2\sum x_t^2,
  \qquad
  \ESS = \sum e_t^2 = \sum y_t^2 - \bh^2\sum x_t^2 .
\end{gathered}
\label{eq:ss}
\end{equation}
Коэффициент детерминации (последнее равенство — тоже только со свободным
членом):
\begin{equation}
  R^2 = \frac{\RSS}{\TSS} = 1 - \frac{\ESS}{\TSS}
      = \frac{\bh^2\sum x_t^2}{\sum y_t^2}.
  \label{eq:r2}
\end{equation}
Свойства оценки наклона и оценка дисперсии ошибок:
\begin{gather}
  \bh = \beta + \frac{\sum x_t\varepsilon_t}{\sum x_t^2},
  \qquad
  \E\bh = \beta,
  \qquad
  \D(\bh) = \frac{\sigma^2}{\sum x_t^2},
  \label{eq:var-b} \\[2pt]
  s^2 = \frac{\sum e_t^2}{n-2},
  \qquad
  s_{\bh}^2 = \frac{s^2}{\sum x_t^2},
  \qquad
  \frac{\bh - \beta}{s_{\bh}} \sim t(n-2).
  \label{eq:s2}
\end{gather}
Тождество \eqref{eq:ss} и равносильность двух формул \eqref{eq:r2} опираются на
\emph{оба} нормальных уравнения МНК, $\sum e_t = 0$ и $\sum e_tX_t = 0$; это
существенно в п.~\ref{sec:noint}.

\begin{remark}
Обозначения $\RSS$ и $\ESS$ не унифицированы. Мы следуем сборнику РАНХиГС:
$\RSS$ (\textit{regression}) — объяснённая сумма квадратов, $\ESS$
(\textit{error}) — остаточная. В ряде книг (например, Johnston, DiNardo, 1997)
те же буквы означают ровно наоборот.
\end{remark}

Понадобятся также свойства дисперсии и ковариации: для неслучайных $c_t$, $d_t$
и попарно некоррелированных $Z_t$
\begin{equation}
  \D\Bigl(\sum c_tZ_t\Bigr) = \sum c_t^2\,\D(Z_t),
  \qquad
  \Cov\Bigl(\sum c_tZ_t,\ \sum d_tZ_t\Bigr) = \sum c_td_t\,\D(Z_t).
  \label{eq:D-props}
\end{equation}

% -------------------------------------------------------------------
\subsection{Регрессия без свободного члена}\label{sec:noint}

В модели $Y_t = \beta X_t + \varepsilon_t$ минимизация $\sum(Y_t - bX_t)^2$ даёт
(вывод — задача~2.8)
\begin{equation}
  \bh = \frac{\sum X_tY_t}{\sum X_t^2},
  \qquad
  \D(\bh) = \frac{\sigma^2}{\sum X_t^2},
  \label{eq:noint}
\end{equation}
а условие первого порядка сводится к единственному нормальному уравнению
$\sum e_tX_t = 0$. Уравнения $\sum e_t = 0$ здесь нет: по свободному члену
дифференцировать нечего. Отсюда три следствия.

\emph{Первое.} Разложение \eqref{eq:ss} перестаёт быть тождеством. Прибавим и
вычтем подогнанное значение $\yh_t = \bh X_t$, чтобы разбить отклонение от
среднего на остаток и отклонение подогнанного значения:
\[
  Y_t - \overline{Y}
  = \bigl(Y_t - \bh X_t\bigr) + \bigl(\bh X_t - \overline{Y}\bigr)
  = e_t + \bigl(\yh_t - \overline{Y}\bigr).
\]
Возведём в квадрат и просуммируем:
\[
  \TSS = \sum\bigl(e_t + (\yh_t - \overline{Y})\bigr)^2
       = \underbrace{\sum e_t^2}_{\ESS}
         + 2\sum e_t\bigl(\yh_t - \overline{Y}\bigr)
         + \underbrace{\sum\bigl(\yh_t - \overline{Y}\bigr)^2}_{\RSS}.
\]
Перекрёстный член распадается на два слагаемых, и первое из них обнуляется
нормальным уравнением, так как $\yh_t = \bh X_t$:
\[
  \sum e_t\bigl(\yh_t - \overline{Y}\bigr)
  = \sum \yh_te_t - \overline{Y}\sum e_t
  = \bh\sum X_te_t - \overline{Y}\sum e_t
  = -\overline{Y}\sum e_t .
\]
Следовательно,
\begin{equation}
  \TSS = \RSS + \ESS - 2\overline{Y}\sum e_t ,
  \label{eq:decomp-noint}
\end{equation}
и разложение верно в точности тогда, когда $\overline{Y}\sum e_t = 0$. В
регрессии со свободным членом перекрёстный член равен нулю целиком:
$\sum \yh_te_t = \ah\sum e_t + \bh\sum X_te_t = 0$ по обоим нормальным
уравнениям, и остаётся \eqref{eq:ss}.

\emph{Второе.} Две формулы \eqref{eq:r2} дают разные числа, и оба могут выходить
за пределы отрезка $[0,1]$: $\RSS/\TSS$ бывает больше единицы, а
$1 - \ESS/\TSS$ — отрицательным. Содержательного коэффициента детерминации в
такой модели попросту нет. Причина в том, что $\TSS$ — это остаточная сумма
прогноза средним, $\yh_t = \overline{Y}$. Модель со свободным членом содержит
такой прогноз: при $\alpha = \overline{Y}$ и $\beta = 0$ она даёт ровно
$\yh_t = \overline{Y}$. МНК выбирает лучшие $\alpha$ и $\beta$, поэтому у такой
модели $\ESS \le \TSS$. Модель без свободного члена этот прогноз не содержит и может оказаться хуже: тогда $\ESS > \TSS$ и
$1 - \ESS/\TSS < 0$.

\emph{Третье.} В модели $Y_t = \beta X_t + \varepsilon_t$ свободный член не
оценивается, а зафиксирован нулём: $\alpha = 0$ и $\yh_t = \bh X_t$. Поэтому
естественная точка отсчёта для неё — не среднее, а ноль: при $\beta = 0$ модель
даёт прогноз $\yh_t = 0$.
Поэтому вместо \emph{центрированных} сумм, то есть отклонений от
$\overline{Y}$, берут \emph{нецентрированные} — отклонения от нуля:
$\sum Y_t^2$, $\sum \yh_t^2$ и $\sum e_t^2$. Для них разложение сохраняется.
По определению остатка $e_t = Y_t - \yh_t$, то есть каждое наблюдение
раскладывается как $Y_t = \yh_t + e_t$. Подставим это в $\sum Y_t^2$ и раскроем
квадрат:
\[
  \sum Y_t^2 = \sum\bigl(\yh_t + e_t\bigr)^2
  = \sum \yh_t^2 + 2\sum \yh_te_t + \sum e_t^2
  = \sum \yh_t^2 + \sum e_t^2 ,
\]
поскольку $\sum \yh_te_t = \bh\sum X_te_t = 0$ по нормальному уравнению.
Отдельные слагаемые $\yh_te_t$ нулю не равны — обнуляется только их сумма. Это
аналог \eqref{eq:ss}, в
котором роль $\TSS$ играет $\sum Y_t^2$, а роль $\RSS$ — $\sum \yh_t^2$.
Отсюда \emph{нецентрированный} коэффициент детерминации
\begin{equation}
  R_u^2 = \frac{\sum \yh_t^2}{\sum Y_t^2} = 1 - \frac{\ESS}{\sum Y_t^2},
  \label{eq:ru2}
\end{equation}
который всегда лежит в $[0,1]$, так как обе суммы в правой части разложения
неотрицательны. Он показывает, насколько модель лучше нулевого прогноза, а не
прогноза средним, поэтому сравнивать его с обычным $R^2$ нельзя. Именно его
печатают эконометрические пакеты для регрессии без константы.

\begin{remark}
Если данные центрированы, то есть $\sum X_t = \sum Y_t = 0$, то
$\sum e_t = \sum Y_t - \bh\sum X_t = 0$ выполняется автоматически, хотя
свободного члена в модели нет. Тогда все формулы п.~\ref{sec:recap} остаются
верными, а $R^2$ совпадает с обычным (задача~2.9, пп.~в, г).
\end{remark}

% -------------------------------------------------------------------
\subsection{Линейные несмещённые оценки и теорема Гаусса--Маркова}\label{sec:gm}

\emph{Линейной} называется оценка вида $\widetilde{\beta} = \sum c_tY_t$ с
неслучайными весами $c_t$. Матожидание суммы равно сумме матожиданий, и
неслучайные $c_t$ выносятся за знак $\E$. По \eqref{eq:model} и предпосылкам
2, 3a $\E Y_t = \E(\alpha + \beta X_t + \varepsilon_t) = \alpha + \beta X_t$,
поэтому
\[
\begin{aligned}
  \E\widetilde{\beta}
  &= \E\Bigl(\sum c_tY_t\Bigr) = \sum c_t\,\E Y_t
   = \sum c_t\bigl(\alpha + \beta X_t\bigr) \\[4pt]
  &= \sum\bigl(\alpha c_t + \beta c_tX_t\bigr)
   = \alpha\sum c_t + \beta\sum c_tX_t .
\end{aligned}
\]
Оценка несмещена, если $\alpha\sum c_t + \beta\sum c_tX_t = \beta$ при любых
$\alpha$ и $\beta$. Это возможно тогда и только тогда, когда коэффициент при
$\alpha$ равен нулю, а при $\beta$ — единице:
\begin{equation}
  \sum c_t = 0, \qquad \sum c_tX_t = \sum c_tx_t = 1
  \label{eq:unb}
\end{equation}
Второе равенство следует из первого: $X_t = x_t + \overline{X}$, а
$\overline{X}$ — константа и выносится за знак суммы, поэтому при
$\sum c_t = 0$
\[
  \sum c_tX_t = \sum c_t\bigl(x_t + \overline{X}\bigr)
  = \sum c_tx_t + \overline{X}\sum c_t
  = \sum c_tx_t .
\]
Для дисперсии заметим, что $Y_t$ некоррелированы и
$\D(Y_t) = \D(\alpha + \beta X_t + \varepsilon_t) = \D(\varepsilon_t) = \sigma^2$
(предпосылки 2, 3a, 3b), поэтому по \eqref{eq:D-props}
\begin{equation}
  \D(\widetilde{\beta}) = \D\Bigl(\sum c_tY_t\Bigr) = \sum c_t^2\,\D(Y_t)
  = \sigma^2\sum c_t^2 .
  \label{eq:var-lin}
\end{equation}
МНК-оценка \eqref{eq:ols} — частный случай с весами
$c_t = x_t\big/\sum x_t^2$: для них условия \eqref{eq:unb} выполнены, а
\eqref{eq:var-lin} даёт уже известную дисперсию $\sigma^2\big/\sum x_t^2$.

\textbf{Теорема Гаусса--Маркова.} При предпосылках 1--3b МНК-оценка обладает
наименьшей дисперсией в классе линейных несмещённых оценок:
\begin{equation}
  \D(\widetilde{\beta}) \ge \frac{\sigma^2}{\sum x_t^2} = \D(\bh),
  \label{eq:gm}
\end{equation}
причём равенство достигается только при $c_t = x_t\big/\sum x_t^2$, то есть
только для МНК-оценки. Такую оценку называют \emph{эффективной}, или BLUE
(\textit{best linear unbiased estimator}).

\emph{Доказательство.} По неравенству Коши--Буняковского и условию \eqref{eq:unb}
\[
  1 = \Bigl(\sum c_tx_t\Bigr)^2 \le \sum c_t^2\,\sum x_t^2
  \quad\Longrightarrow\quad
  \sum c_t^2 \ge \frac{1}{\sum x_t^2},
\]
что вместе с \eqref{eq:var-lin} даёт \eqref{eq:gm}. Равенство в неравенстве
Коши--Буняковского возможно лишь при $c_t = \lambda x_t$, а условие
\eqref{eq:unb} фиксирует $\lambda = 1\big/\sum x_t^2$. $\blacksquare$

\textbf{Объединение независимых оценок.} Пусть $\widehat{\beta}_1$ и
$\widehat{\beta}_2$ — независимые несмещённые оценки $\beta$ с дисперсиями $V_1$
и $V_2$. Оценка
$\widetilde{\beta}_\theta = \theta\widehat{\beta}_1 + (1-\theta)\widehat{\beta}_2$
несмещена при любом $\theta$, а её дисперсия равна
$\theta^2V_1 + (1-\theta)^2V_2$. Приравнивая производную к нулю,
$2\theta V_1 - 2(1-\theta)V_2 = 0$, получаем
\begin{equation}
  \theta^{*} = \frac{V_2}{V_1 + V_2},
  \qquad
  \D\bigl(\widetilde{\beta}^{*}\bigr) = \frac{V_1V_2}{V_1 + V_2}
  = \Bigl(\frac{1}{V_1} + \frac{1}{V_2}\Bigr)^{-1} .
  \label{eq:pool}
\end{equation}
Веса обратно пропорциональны дисперсиям: точнее измеренная оценка получает
больший вес, а дисперсия объединённой оценки меньше каждой из исходных.

% -------------------------------------------------------------------
\subsection{Ковариация оценок и центрирование регрессора}\label{sec:cov}

Обе величины $\overline{Y}$ и $\bh$ — константа плюс линейная комбинация
ошибок. Обозначим $S_{xx} = \sum x_t^2$. Усредняя \eqref{eq:model} по $t$ и
пользуясь \eqref{eq:var-b}, получаем
\[
  \overline{Y} = \alpha + \beta\overline{X} + \sum_t \frac{1}{n}\,\varepsilon_t,
  \qquad
  \bh = \beta + \sum_s \frac{x_s}{S_{xx}}\,\varepsilon_s .
\]
Константы на ковариацию не влияют, а ковариация двух сумм раскрывается по всем
парам $(t, s)$. По предпосылкам 3a, 3b $\Cov(\varepsilon_t, \varepsilon_s)$ равна
$\sigma^2$ при $s = t$ и нулю при $s \neq t$, поэтому в двойной сумме остаются
только слагаемые с $s = t$ — это и есть вторая формула \eqref{eq:D-props}:
\[
\begin{aligned}
  \Cov\bigl(\overline{Y},\, \bh\bigr)
  &= \Cov\Bigl(\sum_t \frac{1}{n}\,\varepsilon_t,\ \sum_s \frac{x_s}{S_{xx}}\,\varepsilon_s\Bigr)
   = \sum_t\sum_s \frac{1}{n}\cdot\frac{x_s}{S_{xx}}\,\Cov(\varepsilon_t, \varepsilon_s) \\[4pt]
  &= \sum_t \frac{1}{n}\cdot\frac{x_t}{S_{xx}}\,\sigma^2
   = \frac{\sigma^2}{n S_{xx}}\sum_t x_t = 0 ,
\end{aligned}
\]
так как сумма отклонений от среднего равна нулю.

Теперь подставим $\ah = \overline{Y} - \overline{X}\,\bh$, где $\overline{X}$ —
константа. Ковариация линейна по каждому аргументу, для дисперсии разности
$\D(U - cV) = \D(U) + c^2\D(V) - 2c\,\Cov(U, V)$, а
$\D(\overline{Y}) = \D\bigl(\sum_t \varepsilon_t/n\bigr) = \sigma^2/n$ по первой
формуле \eqref{eq:D-props}. Поэтому
\[
\begin{aligned}
  \Cov(\ah, \bh)
  &= \Cov\bigl(\overline{Y}, \bh\bigr) - \overline{X}\,\Cov(\bh, \bh)
   = 0 - \overline{X}\,\D(\bh), \\[4pt]
  \D(\ah)
  &= \D(\overline{Y}) + \overline{X}^2\,\D(\bh)
     - 2\overline{X}\,\Cov\bigl(\overline{Y}, \bh\bigr)
   = \frac{\sigma^2}{n} + \overline{X}^2\,\D(\bh) - 0 .
\end{aligned}
\]
Итак,
\begin{equation}
  \Cov(\ah, \bh) = -\overline{X}\,\D(\bh) = -\frac{\overline{X}\,\sigma^2}{\sum x_t^2},
  \qquad
  \D(\ah) = \frac{\sigma^2}{n} + \overline{X}^2\,\D(\bh).
  \label{eq:cov}
\end{equation}
Обе формулы объясняют, зачем центрируют регрессор: при $\overline{X} = 0$ оценки
свободного члена и наклона \emph{некоррелированы}, а дисперсия $\ah$ минимальна.
Сама подгонка при этом не меняется — остатки и $R^2$ у исходной и центрированной
регрессий совпадают (задача~2.21).

% -------------------------------------------------------------------
\subsection{Доверительные интервалы и значимость коэффициентов}\label{sec:ci}

Из \eqref{eq:s2} при предпосылке 3c доверительный интервал уровня 95\,\% имеет
вид
\begin{equation}
  \bh \pm t_{\text{кр}}\,s_{\bh},
  \qquad
  t_{\text{кр}} = t_{0{,}025}(n-2);
  \label{eq:ci}
\end{equation}
для $\alpha$ он строится так же с $s_{\ah}$. Коэффициент называют
\emph{значимым} на уровне 5\,\%, если отвергается гипотеза $H_0\colon \beta = 0$,
то есть $|t| = |\bh/s_{\bh}| > t_{\text{кр}}$. Три формулировки равносильны:
$|t| > t_{\text{кр}}$; $P$-значение $P = \Pr\bigl\{|t(n-2)| > |\bh/s_{\bh}|\bigr\}$
меньше 0,05; интервал \eqref{eq:ci} не содержит нуля.

% -------------------------------------------------------------------
\subsection{Другие критерии подгонки}\label{sec:lad}

МНК минимизирует $S(a,b) = \sum(Y_t - a - bX_t)^2$ — строго выпуклую функцию
(семинар~1), поэтому при условии, что не все $X_t$ равны, решение существует и
единственно. Метод наименьших модулей (МНМ) минимизирует
\begin{equation}
  F(a,b) = \sum\bigl|Y_t - a - bX_t\bigr| .
  \label{eq:lad}
\end{equation}
Эта функция выпукла, но кусочно-линейна и \emph{не строго} выпукла: там, где ни
одно слагаемое не меняет знак, она линейна по параметрам. Поэтому минимум может
достигаться не в одной точке, а на целом множестве (задача~2.11).
\fi

% ===================================================================
\section{Задачи}

% -------------------------------------------------------------------
\problem{2.8}

Выведите непосредственно формулу для оценки коэффициента наклона в регрессии
\textbf{без свободного члена}, т.\,е. найдите оценку параметра $\beta$ в регрессии
\[
  Y_t = \beta X_t + \varepsilon_t
\]
минимизацией суммы квадратов отклонений $\sum \bigl(Y_t - \yh_t\bigr)^2$.

\begin{solution}
Минимизируем
\[
  S(b) = \sum\bigl(Y_t - bX_t\bigr)^2 .
\]
Условие первого порядка
\[
  S'(b) = -2\sum\bigl(Y_t - bX_t\bigr)X_t = 0
\]
— это и есть единственное нормальное уравнение $\sum e_tX_t = 0$. Решая его
относительно $b$, получаем МНК-оценку
\[
  \boxed{\ \bh = \frac{\sum X_tY_t}{\sum X_t^2}\ }
\]
Это минимум, так как $S''(b) = 2\sum X_t^2 > 0$.

\medskip
Свойства оценки выводятся так же, как в семинаре~1, только с $X_t$ вместо
$x_t$. Подставим $Y_t = \beta X_t + \varepsilon_t$:
\[
  \bh = \frac{\sum X_t\bigl(\beta X_t + \varepsilon_t\bigr)}{\sum X_t^2}
      = \frac{\beta\sum X_t^2 + \sum X_t\varepsilon_t}{\sum X_t^2}
      = \beta + \frac{\sum X_t\varepsilon_t}{\sum X_t^2}.
\]
Веса $X_t/\sum X_t^2$ неслучайны (предпосылка~2), поэтому по предпосылке~3a
$\E\bh = \beta$. В дисперсии константа $\beta$ пропадает, а по
\eqref{eq:D-props} и предпосылкам 3a, 3b множитель $1/\sum X_t^2$ выносится в
квадрате и дисперсия суммы равна сумме дисперсий:
\[
\begin{aligned}
  \D(\bh) &= \D\Bigl(\beta + \frac{\sum X_t\varepsilon_t}{\sum X_t^2}\Bigr)
           = \frac{1}{\bigl(\sum X_t^2\bigr)^2}\,\D\Bigl(\sum X_t\varepsilon_t\Bigr)
           = \frac{1}{\bigl(\sum X_t^2\bigr)^2}\sum X_t^2\,\D(\varepsilon_t) \\[4pt]
          &= \frac{\sigma^2\sum X_t^2}{\bigl(\sum X_t^2\bigr)^2}
           = \frac{\sigma^2}{\sum X_t^2}.
\end{aligned}
\]
Это и есть формулы \eqref{eq:noint}.

\medskip
Второго нормального уравнения, $\sum e_t = 0$, здесь нет: по свободному члену
дифференцировать нечего. Все следствия этого разобраны в п.~\ref{sec:noint} —
разложение \eqref{eq:ss} нарушается на величину $2\overline{Y}\sum e_t$, две
формулы \eqref{eq:r2} перестают быть эквивалентными, и осмысленным остаётся
только нецентрированный коэффициент \eqref{eq:ru2}. Численно это видно в
следующей задаче.
\end{solution}

% -------------------------------------------------------------------
\problem{2.9}

Для наблюдений

\begin{center}
\begin{tabular}{l *{10}{c}}
\toprule
$t$ & 1 & 2 & 3 & 4 & 5 & 6 & 7 & 8 & 9 & 10 \\
\midrule
$Y_t$ & 70 & 65 & 55 & 60 & 50 & 35 & 40 & 30 & 25 & 32 \\
$X_t$ &  5 & 11 & 15 & 17 & 20 & 22 & 25 & 27 & 30 & 35 \\
\bottomrule
\end{tabular}
\end{center}

\noindent вычислите следующие величины:
\begin{enumerate}[label=\asbuk*), leftmargin=2.2em, itemsep=3pt]
  \item коэффициент детерминации $R^2$ в регрессии $Y_t$ на $X_t$
        \textbf{при наличии} свободного члена;
  \item коэффициент детерминации $R^2$ в регрессии $Y_t$ на $X_t$
        \textbf{при отсутствии} свободного члена;
  \item коэффициент детерминации $R^2$ в регрессии $y_t$ на $x_t$
        \textbf{при наличии} свободного члена, где $y_t$ и $x_t$ --- отклонения
        переменных $Y_t$ и $X_t$ от их средних значений;
  \item коэффициент детерминации $R^2$ в регрессии $y_t$ на $x_t$
        \textbf{при отсутствии} свободного члена.
\end{enumerate}

\begin{excel}
Надстройка подключается один раз: \ui{Файл} $\to$ \ui{Параметры} $\to$
\ui{Надстройки} $\to$ \ui{Пакет анализа}.

\smallskip
Занесите данные двумя столбцами. \textbf{Внимательно с порядком:} в таблице
второй столбец --- это $Y$ (70, 65, \dots), третий --- $X$ (5, 11, \dots),
а не наоборот.
\begin{enumerate}[label=\asbuk*), leftmargin=2.2em, itemsep=2pt, topsep=3pt]
  \item \ui{Данные} $\to$ \ui{Анализ данных} $\to$ \ui{Регрессия}.
        \ui{Входной интервал Y} --- столбец $Y$, \ui{Входной интервал X} ---
        столбец $X$. В отчёте: \ui{Y-пересечение} $=79{,}95$, коэффициент при
        $X$ $=-1{,}63$, \ui{R-квадрат} $=0{,}8607$.
  \item Тот же диалог, но включите флажок \ui{Константа --- ноль}.
        Коэффициент станет $1{,}66$.
  \item Добавьте два столбца отклонений и повторите п.~а) для них:\newline
        \xl{=B2-СРЗНАЧ(\$B\$2:\$B\$11)} \quad и \quad \xl{=C2-СРЗНАЧ(\$C\$2:\$C\$11)}
  \item То же, что в), но с флажком \ui{Константа --- ноль}.
\end{enumerate}

\smallskip
Сравните \ui{R-квадрат} из четырёх отчётов между собой и с ответами ниже.
В пункте б) Excel покажет \textbf{не} то, что даёт формула
$R^2 = \RSS/\TSS$, --- разберитесь, почему.
\end{excel}

\begin{solution}
\textbf{Подготовительные суммы} ($n = 10$):
\[
  \textstyle\sum X_t = 207,\quad \sum Y_t = 462,\quad
  \sum X_t^2 = 5023,\quad \sum Y_t^2 = 23\,624,\quad \sum X_tY_t = 8360,
\]
откуда $\overline{X} = 20{,}7$, $\overline{Y} = 46{,}2$ и, по
\eqref{eq:dev-sums},
\[
\begin{aligned}
  \sum x_t^2  &= 5023 - \frac{207^2}{10} = 738{,}1, \\[2pt]
  \sum x_ty_t &= 8360 - \frac{207\cdot 462}{10} = -1203{,}4, \\[2pt]
  \sum y_t^2  &= 23\,624 - \frac{462^2}{10} = 2279{,}6 .
\end{aligned}
\]

\medskip
\textbf{а)} По \eqref{eq:ols}, второй строке \eqref{eq:ss} и \eqref{eq:r2}
\[
  \bh = \frac{-1203{,}4}{738{,}1} = -1{,}63,
  \qquad
  \ah = 46{,}2 + 1{,}63\cdot 20{,}7 = 79{,}95,
  \qquad
  \yh_t = 79{,}95 - 1{,}63\,X_t,
\]
\[
  \RSS = \bh^2\sum x_t^2 = 1962{,}0,
  \qquad
  \TSS = 2279{,}6,
  \qquad
  R^2 = \frac{1962{,}0}{2279{,}6} = 0{,}8607 .
\]

\medskip
\textbf{б)} Теперь модель $Y_t = \beta X_t + \varepsilon_t$, и по
\eqref{eq:noint}
\[
  \bh = \frac{8360}{5023} = 1{,}66,
  \qquad
  \yh_t = 1{,}66\,X_t .
\]
Суммы квадратов считаем по определениям из первой строки \eqref{eq:ss}:
формулы через отклонения из второй строки здесь неприменимы (например,
$\bh^2\sum x_t^2 = 2044{,}6$ — это не $\RSS$ этой модели). Получаем
\[
  \RSS = \sum\bigl(\yh_t - \overline{Y}\bigr)^2 = 3424{,}7,
  \qquad
  \ESS = \sum e_t^2 = 9710{,}1,
  \qquad
  \sum e_t = 117{,}5 \neq 0 .
\]
Сумма $\RSS + \ESS = 13\,134{,}8$ не равна $\TSS = 2279{,}6$; расхождение в
точности описывается формулой \eqref{eq:decomp-noint}:
$2\overline{Y}\sum e_t = 2\cdot 46{,}2\cdot 117{,}5 = 10\,855{,}2$. Поэтому две
формулы \eqref{eq:r2} дают разные и бессмысленные значения:
\[
  \frac{\RSS}{\TSS} = \frac{3424{,}7}{2279{,}6} = 1{,}50 > 1,
  \qquad
  1 - \frac{\ESS}{\TSS} = 1 - \frac{9710{,}1}{2279{,}6} = -3{,}26 < 0 .
\]
Корректной характеристикой здесь служит нецентрированный коэффициент
\eqref{eq:ru2}:
\[
  R_u^2 = 1 - \frac{9710{,}1}{23\,624} = 0{,}5890 ,
\]
и именно это число печатает Excel с флажком \ui{Константа --- ноль}: ни
$1{,}50$, ни $-3{,}26$. Пакет молча подменяет $\TSS$ на $\sum Y_t^2$; сравнивать
такой $R^2$ с $R^2$ из пункта~а) нельзя.

\medskip
\textbf{в), г)} Для центрированных данных $\overline{x} = \overline{y} = 0$,
поэтому по \eqref{eq:ols} $\widehat{\alpha}' = \overline{y} - \bh\overline{x} = 0$:
свободный член оценивается точно нулём, и пункты в) и г) описывают одну и ту же
прямую $\widehat{y}_t = \bh x_t$. По замечанию к п.~\ref{sec:noint} равенство
$\sum e_t = 0$ здесь выполняется и без константы, так что \eqref{eq:ss} и
\eqref{eq:r2} остаются в силе. Наклон и суммы квадратов те же, что в п.~а):
\[
  \bh = \frac{\sum x_ty_t}{\sum x_t^2} = -1{,}63,
  \qquad
  R^2 = \frac{\bh^2\sum x_t^2}{\sum y_t^2} = 0{,}8607 .
\]

\medskip
\textbf{Вывод.}
\begin{center}
\begin{tabular}{clcl}
\toprule
Пункт & Регрессия & $\bh$ & $R^2$ \\
\midrule
а) & $Y$ на $X$ со свободным членом  & $-1{,}63$ & $0{,}8607$ \\
б) & $Y$ на $X$ без свободного члена & $\phantom{-}1{,}66$ & $1{,}50$ \ \ или \ $-3{,}26$ \\
в) & $y$ на $x$ со свободным членом  & $-1{,}63$ & $0{,}8607$ \\
г) & $y$ на $x$ без свободного члена & $-1{,}63$ & $0{,}8607$ \\
\bottomrule
\end{tabular}
\end{center}
Центрирование делает свободный член излишним: он и так равен нулю. А отбрасывание
константы у нецентрированных данных (пункт~б) меняет модель по существу и
разрушает само понятие $R^2$.
\end{solution}

% -------------------------------------------------------------------
\problem{2.16 (кривая Филлипса)}

Так называемая кривая Филлипса описывает связь темпа роста зарплаты и уровня
безработицы. А именно,
\[
  \delta w_t = \beta_1 + \beta_2 \frac{1}{u_t} + \varepsilon_t,
\]
где $w_t$ --- уровень заработной платы,
$\delta w_t = 100\,(w_t - w_{t-1})/w_{t-1}$ --- темп роста зарплаты (в процентах),
$u_t$ --- процент безработных в год $t$. Теория предполагает, что
$\beta_1 < 0$ и $\beta_2 > 0$.

\begin{center}
\small
\begin{tabular}{l *{9}{c}}
\toprule
$t$   & 1 & 2 & 3 & 4 & 5 & 6 & 7 & 8 & 9 \\
\midrule
$w_t$ & 1,62 & 1,65 & 1,79 & 1,94 & 2,03 & 2,12 & 2,26 & 2,44 & 2,57 \\
$u_t$ & 1,0 & 1,4 & 1,1 & 1,5 & 1,5 & 1,2 & 1,0 & 1,1 & 1,3 \\
\midrule
$t$   & 10 & 11 & 12 & 13 & 14 & 15 & 16 & 17 & 18 \\
\midrule
$w_t$ & 2,66 & 2,73 & 2,80 & 2,92 & 3,02 & 3,13 & 3,28 & 3,43 & 3,58 \\
$u_t$ & 1,8 & 1,9 & 1,5 & 1,4 & 1,8 & 1,1 & 1,5 & 1,3 & 1,4 \\
\bottomrule
\end{tabular}
\end{center}

\noindent Используя данные для некоторой страны:
\begin{enumerate}[label=\asbuk*), leftmargin=2.2em, itemsep=3pt]
  \item Найдите оценки коэффициентов уравнения и проверьте наличие значимой
        связи между $\delta w$ и $u$.
  \item Найдите <<естественный уровень безработицы>>, т.\,е. такой уровень
        безработицы, при котором $\delta w = 0$.
  \item Когда изменения в уровне безработицы оказывали наибольшее (наименьшее)
        влияние на темп изменения зарплаты?
  \item Найдите 95\,\%-доверительные интервалы для $\beta_1$ и $\beta_2$.
\end{enumerate}

\begin{excel}
Столбцы $w_t$ и $u_t$ вводятся как есть, ещё два считаются формулами
(первая строка остаётся пустой --- темп роста для $t=1$ не определён):
\[
  \text{\xl{=100*(A3-A2)/A2}}
  \qquad\text{и}\qquad
  \text{\xl{=1/B3}}
\]
\textbf{Множитель 100 и вид регрессора важны.} Если считать $\delta w$ долей
(без \xl{100*}) или брать $100/u$ вместо $1/u$, то $R^2$, $t$-статистики и
$P$-значения останутся верными --- они не зависят от масштаба, --- а вот
коэффициенты выйдут в других единицах и с ответом не сойдутся.

\smallskip
Дальше \ui{Данные} $\to$ \ui{Анализ данных} $\to$ \ui{Регрессия} по 17 строкам,
с флажком \ui{Остатки}. В отчёте сразу есть всё для пунктов а) и г):
коэффициенты, стандартные ошибки, $P$-значения и границы
\ui{Нижние 95\%} / \ui{Верхние 95\%}.

\smallskip
\textbf{Проверьте отчёт вручную.} Стандартную ошибку наклона
\[
  s_{\bh} = \sqrt{\frac{\ESS/(n-2)}{\sum (x_t-\overline{x})^2}}
\]
можно собрать из тех же остатков в три шага:
\begin{center}\small
\begin{tabular}{@{}llr@{}}
\toprule
Величина & Формула & Значение \\
\midrule
$s^2 = \ESS/(n-2)$              & \xl{=СУММКВ(остатки)/15}   & 3,0098 \\
$\sum (x_t-\overline{x})^2$     & \xl{=ДИСП.Г(столбец 1/u)*17} & 0,3009 \\
$s_{\bh}$                       & \xl{=КОРЕНЬ(первое/второе)} & 3,1626 \\
\bottomrule
\end{tabular}
\end{center}
Последнее число --- ровно то, что Excel напечатал в графе
\ui{Стандартная ошибка}. Здесь \xl{ДИСП.Г}$\,\cdot n$ --- это и есть
$\sum (x_t-\overline{x})^2$, поскольку \xl{ДИСП.Г} делит на $n$, а не на $n-1$.
\end{excel}

% Решение начинается с новой страницы: иначе таблица регрессии попадает в
% конец страницы внутри рамки решения и наезжает на колонтитул.
\ifsolutions\clearpage\fi
\begin{solution}
Темп роста $\delta w_t$ определён начиная с $t = 2$, поэтому в регрессии
участвуют $n = 17$ наблюдений, а число степеней свободы равно $n - 2 = 15$.
Регрессор — величина $1/u_t$; модель линейна по параметрам, поэтому все формулы
п.~\ref{sec:recap} применимы без изменений.

\medskip
\textbf{а)} Оценка параметров даёт следующий результат.

\begin{center}
\begin{tabular}{lrrrr}
\toprule
\multicolumn{5}{l}{\itshape Зависимая переменная: $\delta w$;\quad $n = 17$} \\
\midrule
Переменная & Коэффициент & Ст. ошибка & $t$-статистика & $P$-значение \\
\midrule
const & $-1{,}034$ & 2,3716 & $-0{,}4362$ & 0,6689 \\
$1/u$ & $\phantom{-}7{,}896$ & 3,1626 & $\phantom{-}2{,}4966$ & 0,0247 \\
\midrule
\multicolumn{5}{l}{$R^2 = 0{,}294$;\quad $\ESS = 45{,}15$;\quad $\TSS = 63{,}91$;\quad $F(1,15) = 6{,}23$} \\
\bottomrule
\end{tabular}
\end{center}

\noindent Итоговое уравнение:
$\widehat{\delta w} = -1{,}034 + 7{,}896/u$.

По критерию п.~\ref{sec:ci} коэффициент при $1/u$ значим
($P = 0{,}0247 < 0{,}05$), а свободный член от нуля достоверно не отличается
($P = 0{,}6689$). Связь между $\delta w$ и $u$ есть: с ростом безработицы темп
роста зарплаты падает. Знаки обоих коэффициентов согласуются с теорией
($\widehat{\beta}_1 < 0$, $\widehat{\beta}_2 > 0$).

\medskip
\textbf{б)} «Естественный уровень безработицы» — корень уравнения
$\widehat{\delta w} = 0$:
\[
  -1{,}034 + 7{,}896\,\frac{1}{u} = 0
  \quad\Longrightarrow\quad
  \boxed{\ u_0 = \frac{7{,}896}{1{,}034} = 7{,}63\,\%\ }
\]

\medskip
\textbf{в)} Модель нелинейна по $u$, поэтому эффект зависит от уровня
безработицы:
\[
  \frac{\partial\,\delta w}{\partial u} = -\frac{7{,}896}{u^2} .
\]
По модулю он тем больше, чем меньше $u$. Наибольшее влияние — в годы с
наименьшей безработицей ($u = 1{,}0$: 1-й и 7-й годы,
$\partial\,\delta w/\partial u = -7{,}90$), наименьшее — в год с наибольшей
($u = 1{,}9$: 11-й год, $-2{,}19$). Первый год при этом выпадает из выборки:
для него $\delta w$ не определён.

\medskip
\textbf{г)} По \eqref{eq:ci} при 15 степенях свободы $t_{0{,}025}(15) = 2{,}131$:
\[
\begin{aligned}
  \beta_1 &= -1{,}034 \pm 2{,}131 \cdot 2{,}372 = -1{,}034 \pm 5{,}055
          &&\Longrightarrow\quad [-6{,}09;\ \ 4{,}02], \\[2pt]
  \beta_2 &= \phantom{-}7{,}896 \pm 2{,}131 \cdot 3{,}163 = 7{,}896 \pm 6{,}741
          &&\Longrightarrow\quad [\phantom{-}1{,}15;\ 14{,}64].
\end{aligned}
\]
Первый интервал накрывает ноль, второй — нет; это согласуется с $P$-значениями
из пункта~а) и с критерием значимости п.~\ref{sec:ci}.

\begin{remark}
В решебнике в этом пункте взято $t_{0{,}025}(16) = 2{,}12$, что даёт
$-1{,}03 \pm 5{,}03$ и $7{,}90 \pm 6{,}70$. Это опечатка в числе степеней свободы:
наблюдений 17 (первое теряется при переходе к темпу роста), оцениваемых
параметров 2, значит $\mathrm{df} = 15$, а не 16. На выводы расхождение не влияет.
\end{remark}
\end{solution}

% ===================================================================
\section{Дополнительные задачи}

Задачи этого раздела не входят в основной ход семинара: они разбираются, если
остаётся время, или предлагаются для самостоятельной работы.

% -------------------------------------------------------------------
\problem{2.23}

Два исследователя, работая независимо друг от друга, изучают одну и ту же
регрессионную модель
\[
  Y_t = \alpha + \beta X_t + \varepsilon_t,
\]
для которой выполнены все условия классической модели. В таблице приведены
результаты, полученные ими на основе независимых выборок.

\begin{center}
\begin{tabular}{cc}
\toprule
Выборка I & Выборка II \\
\midrule
$n = 20$              & $n = 20$ \\
$\sum X_t = 100$      & $\sum X_t = 200$ \\
$\sum X_t^2 = 600$    & $\sum X_t^2 = 2400$ \\
$\sum Y_t = 500$      & $\sum Y_t = 500$ \\
$\widehat{\beta}_{\mathrm{I}} = 2$ & $\widehat{\beta}_{\mathrm{II}} = 2{,}5$ \\
\bottomrule
\end{tabular}
\end{center}

Узнав о работе друг друга, они решают вывести единую оценку параметра $\beta$.
Первый исследователь предлагает взять
\[
  \widetilde{\beta} = \tfrac{1}{2}\bigl(\widehat{\beta}_{\mathrm{I}} + \widehat{\beta}_{\mathrm{II}}\bigr).
\]
Второй исследователь считает, что весовые коэффициенты первой и второй оценок
выбраны неэффективно, и можно построить несмещённую оценку с меньшей дисперсией.
Научный руководитель этих исследователей утверждает, что он знает способ ещё
улучшить общую оценку.
\begin{enumerate}[label=\asbuk*), leftmargin=2.2em, itemsep=3pt]
  \item Какую оценку предлагает использовать второй исследователь?
  \item Какую оценку предлагает использовать научный руководитель?
\end{enumerate}
Оцените улучшение точности оценок пп.~а), б) по сравнению с оценкой
$\widetilde{\beta}$.

\begin{solution}
Исследователи используют МНК-оценки, поэтому по \eqref{eq:dev-sums} и
\eqref{eq:var-b}
\[
  \sum_{\mathrm{I}} x_t^2 = 600 - \frac{100^2}{20} = 100,
  \qquad
  \sum_{\mathrm{II}} x_t^2 = 2400 - \frac{200^2}{20} = 400,
\]
\[
  V_1 = \D\bigl(\widehat{\beta}_{\mathrm{I}}\bigr) = \frac{\sigma^2}{100},
  \qquad
  V_2 = \D\bigl(\widehat{\beta}_{\mathrm{II}}\bigr) = \frac{\sigma^2}{400}.
\]

\medskip
\textbf{а)} Второй исследователь предлагает взвесить оценки оптимально. По
\eqref{eq:pool}
\[
  \theta^{*} = \frac{V_2}{V_1 + V_2} = \frac{\sigma^2/400}{\sigma^2/100 + \sigma^2/400}
             = \frac{1}{5},
  \qquad
  \boxed{\ \widetilde{\beta}^{*} = \tfrac{1}{5}\cdot 2 + \tfrac{4}{5}\cdot 2{,}5 = 2{,}4\ }
\]
Вторая выборка даёт вчетверо больший разброс регрессора, её оценка вчетверо
точнее и получает вес $4/5$. Дисперсия по \eqref{eq:pool} равна
$V_1V_2/(V_1+V_2) = \sigma^2/500 = 0{,}002\,\sigma^2$.

\medskip
\textbf{б)} Научный руководитель предлагает объединить не оценки, а \emph{данные}:
по теореме Гаусса--Маркова \eqref{eq:gm} эффективна МНК-оценка по всем
$n = 40$ наблюдениям. Недостающие суммы $\sum X_tY_t$ восстанавливаются из
\eqref{eq:ols} и \eqref{eq:dev-sums}: $\sum X_tY_t = \bh\sum x_t^2 + \frac{1}{n}\sum X_t\sum Y_t$,
то есть
\[
  \sum_{\mathrm{I}} X_tY_t = 2\cdot 100 + \frac{100\cdot 500}{20} = 2700,
  \qquad
  \sum_{\mathrm{II}} X_tY_t = 2{,}5\cdot 400 + \frac{200\cdot 500}{20} = 6000 .
\]
Для объединённой выборки $n = 40$, $\sum X_t = 300$, $\sum Y_t = 1000$,
$\sum X_t^2 = 3000$, $\sum X_tY_t = 8700$, поэтому
\[
  \sum x_t^2 = 3000 - \frac{300^2}{40} = 750,
  \qquad
  \sum x_ty_t = 8700 - \frac{300\cdot 1000}{40} = 1200,
\]
\[
  \boxed{\ \bh = \frac{1200}{750} = 1{,}6\ }
\]
и $\D(\bh) = \sigma^2/750 \approx 0{,}001333\,\sigma^2$.

\medskip
\textbf{Сравнение точности.} Для оценки с равными весами по \eqref{eq:var-lin}
$\D(\widetilde{\beta}) = \tfrac14 V_1 + \tfrac14 V_2 = 0{,}003125\,\sigma^2$.
\begin{center}
\begin{tabular}{lcc}
\toprule
Оценка & Дисперсия & Выигрыш против $\widetilde{\beta}$ \\
\midrule
$\widetilde{\beta} = 2{,}25$ (равные веса) & $0{,}003125\,\sigma^2$ & --- \\
$\widetilde{\beta}^{*} = 2{,}4$ (оптимальные веса) & $0{,}002000\,\sigma^2$ & $-36\,\%$ \\
$\bh = 1{,}6$ (объединённая выборка) & $0{,}001333\,\sigma^2$ & $-57\,\%$ \\
\bottomrule
\end{tabular}
\end{center}

\begin{remark}
Оценка $\bh = 1{,}6$ лежит вне отрезка между 2 и 2,5, то есть не является
никакой выпуклой комбинацией исходных оценок. Объединение \emph{данных} — это не
усреднение \emph{оценок}: объединённая регрессия использует ещё и различие
средних $\overline{X}$ между выборками, которого взвешивание оценок «не видит».
\end{remark}
\end{solution}

% -------------------------------------------------------------------
\problem{2.10}

Предположим, что модель
\[
  Y_t = \alpha + \beta X_t + \varepsilon_t, \qquad t = 1,\dots,n
\]
удовлетворяет условиям классической регрессии. Рассматривается следующая оценка
коэффициента $\beta$:
\[
  \widetilde{\beta} = \frac{1}{n}\sum_{t=1}^{n} \frac{Y_t - \overline{Y}}{X_t - \overline{X}} .
\]
\begin{enumerate}[label=\asbuk*), leftmargin=2.2em, itemsep=3pt]
  \item Является ли оценка $\widetilde{\beta}$ несмещённой? Является ли она линейной?
  \item Вычислите дисперсию оценки $\widetilde{\beta}$.
  \item Проверьте теорему Гаусса--Маркова, сравнив полученную дисперсию оценки
        $\widetilde{\beta}$ с дисперсией МНК-оценки
        $\sigma^2\big/\sum_{t=1}^{n}(X_t - \overline{X})^2$.
\end{enumerate}

\begin{solution}
Всюду предполагается $x_t = X_t - \overline{X} \neq 0$, иначе оценка не
определена.

\medskip
\textbf{а)} Раскроем $\overline{Y} = \frac{1}{n}\sum_s Y_s$ и соберём
коэффициенты при $Y_t$:
\[
  \widetilde{\beta}
  = \frac{1}{n}\sum_t \frac{Y_t}{x_t}
    - \frac{1}{n^2}\Bigl(\sum_s Y_s\Bigr)\Bigl(\sum_t \frac{1}{x_t}\Bigr)
  = \sum_t c_tY_t,
  \qquad
  c_t = \frac{1}{n}\left(\frac{1}{x_t} - \frac{1}{n}\sum_s \frac{1}{x_s}\right).
\]
Веса $c_t$ неслучайны (предпосылка~2), значит оценка \textbf{линейна}. Проверим
условия несмещённости \eqref{eq:unb}:
\[
  \sum_t c_t = \frac{1}{n}\left(\sum_t\frac{1}{x_t} - \sum_s\frac{1}{x_s}\right) = 0,
  \qquad
  \sum_t c_tx_t
  = \frac{1}{n}\left(n - \frac{1}{n}\Bigl(\sum_s\frac{1}{x_s}\Bigr)\sum_t x_t\right)
  = 1,
\]
поскольку $\sum_t x_t = 0$. Оба условия выполнены, поэтому оценка
\textbf{несмещённая}: $\E\widetilde{\beta} = \beta$.

\medskip
\textbf{б)} По \eqref{eq:var-lin}
\[
  \D(\widetilde{\beta}) = \sigma^2\sum_t c_t^2
  = \frac{\sigma^2}{n^2}\sum_{t=1}^{n}
    \left(\frac{1}{x_t} - \frac{1}{n}\sum_{s=1}^{n}\frac{1}{x_s}\right)^{\!2}.
\]

\medskip
\textbf{в)} Оценка линейна и несмещённа, поэтому к ней применима теорема
Гаусса--Маркова \eqref{eq:gm}:
\[
  \D(\widetilde{\beta}) = \sigma^2\sum_t c_t^2 \ \ge\ \frac{\sigma^2}{\sum_t x_t^2}
  = \D(\bh).
\]
Равенство было бы возможно только при $c_t \propto x_t$, тогда как здесь
$c_t$ пропорциональны $1/x_t$ с точностью до сдвига. Значит, МНК-оценка строго
эффективнее (кроме вырожденных наборов данных), и теорема подтверждается.
\end{solution}

% -------------------------------------------------------------------
\problem{2.11}

Приведите пример набора данных $(X_t, Y_t)$, для которого решение задачи поиска
параметров $\alpha$, $\beta$, минимизирующих функционал
\[
  F = \sum_{t=1}^{n} \bigl|Y_t - (\alpha + \beta X_t)\bigr|,
\]
не единственно.

\begin{solution}
Возьмём $n = 3$ и вершины равностороннего треугольника со стороной 2:
$(X_1,Y_1) = (0,0)$, $(X_2,Y_2) = (0,2)$, $(X_3,Y_3) = (\sqrt{3},1)$. Тогда
функционал \eqref{eq:lad} равен
\[
  F = |\alpha| + |2 - \alpha| + \bigl|1 - \alpha - \sqrt{3}\beta\bigr| .
\]
Первые два слагаемых оцениваются снизу как
$|\alpha| + |2-\alpha| \ge 2$, причём равенство достигается при любом
$\alpha \in [0,2]$; третье слагаемое неотрицательно и обращается в нуль при
$\alpha + \sqrt{3}\beta = 1$. Следовательно, $F_{\min} = 2$ и достигается на
всём множестве
\[
  \alpha \in [0,2], \qquad \beta = \frac{1-\alpha}{\sqrt{3}} ,
\]
то есть на всех прямых, проходящих через точку $(\sqrt{3},1)$ и пересекающих ось
$Y$ между 0 и 2. Например, $\alpha = 0$, $\beta = 1/\sqrt{3}$ и $\alpha = 2$,
$\beta = -1/\sqrt{3}$. Решений \textbf{бесконечно много}: сумма модулей «плоская»
на отрезке, и это в точности та неединственность, о которой сказано в
п.~\ref{sec:lad}.
\end{solution}

% -------------------------------------------------------------------
\problem{2.21}

Проведены две регрессии:
\[
  Y_t = \alpha + \beta X_t + \varepsilon_t
  \qquad\text{и}\qquad
  Y_t = \alpha' + \beta' x_t + \varepsilon'_t, \qquad t = 1,\dots,T,
\]
где $x_t = X_t - \overline{X}$.
\begin{enumerate}[label=\asbuk*), leftmargin=2.2em, itemsep=3pt]
  \item По известным МНК-оценкам $\ah$, $\bh$ параметров $\alpha$, $\beta$
        в первой регрессии найдите МНК-оценки $\widehat{\alpha}'$,
        $\widehat{\beta}'$ параметров $\alpha'$, $\beta'$ во второй регрессии.
  \item Найдите $\Cov(\widehat{\alpha}', \widehat{\beta}')$.
\end{enumerate}

\begin{solution}
\textbf{а)} Регрессор второй регрессии центрирован: $\overline{x} = 0$. По
\eqref{eq:ols}
\[
  \widehat{\beta}'
  = \frac{\sum (x_t - \overline{x})(Y_t - \overline{Y})}{\sum (x_t - \overline{x})^2}
  = \frac{\sum x_t y_t}{\sum x_t^2} = \bh,
\]
\[
  \widehat{\alpha}' = \overline{Y} - \widehat{\beta}'\,\overline{x} = \overline{Y}
  = \bigl(\overline{Y} - \bh\,\overline{X}\bigr) + \bh\,\overline{X}
  = \ah + \bh\,\overline{X} .
\]
Наклон не меняется, а свободный член становится равен $\overline{Y}$: точка
отсчёта переносится в $\overline{X}$.

\medskip
\textbf{б)} По формуле \eqref{eq:cov}, применённой ко второй регрессии, у которой
среднее регрессора равно нулю,
\[
  \Cov(\widehat{\alpha}', \widehat{\beta}')
  = -\overline{x}\,\D(\widehat{\beta}') = 0 .
\]
Это и есть причина, по которой регрессоры центрируют (п.~\ref{sec:cov}).
\end{solution}

% -------------------------------------------------------------------
\problem{2.13}

Рассмотрим модель регрессии без константы
\[
  Y_t = \beta X_t + \varepsilon_t, \qquad t = 1,\dots,n .
\]
\begin{enumerate}[label=\asbuk*), leftmargin=2.2em, itemsep=3pt]
  \item Найдите оценки метода наименьших квадратов для $\beta$ и $\sigma^2$.
  \item Найдите дисперсию оценки $\bh$.
  \item Покажите, что статистика $\dfrac{\bh - \beta}{s_{\bh}}$ имеет
        распределение $t(n-1)$.
  \item Приведите примеры данных, для которых: значение коэффициента $R^2$,
        рассчитанное по формуле $R^2 = \RSS/\TSS$, отличается от значения $R^2$,
        рассчитанного по формуле $R^2 = 1 - \ESS/\TSS$; значение $R^2$,
        рассчитанное по формуле $R^2 = \RSS/\TSS$, больше 1; значение $R^2$,
        рассчитанное по формуле $R^2 = 1 - \ESS/\TSS$, меньше 0.
\end{enumerate}

\begin{solution}
\textbf{а)} Оценка наклона дана формулой \eqref{eq:noint}:
$\bh = \sum X_tY_t \big/ \sum X_t^2$. Для оценки $\sigma^2$ выразим сумму
квадратов остатков, пользуясь нормальным уравнением
$\sum X_tY_t = \bh\sum X_t^2$:
\[
  \sum e_t^2 = \sum\bigl(Y_t - \bh X_t\bigr)^2
  = \sum Y_t^2 - 2\bh\sum X_tY_t + \bh^2\sum X_t^2
  = \sum Y_t^2 - \bh^2\sum X_t^2 .
\]
Оба слагаемых усредняем по \eqref{eq:D-props} и правилу
$\E\bigl[Z^2\bigr] = \D(Z) + (\E Z)^2$:
\[
  \E\sum Y_t^2 = \sum\Bigl(\D(Y_t) + \bigl(\E Y_t\bigr)^2\Bigr)
               = n\sigma^2 + \beta^2\sum X_t^2,
  \qquad
  \E\bigl[\bh^2\bigr]\sum X_t^2 = \sigma^2 + \beta^2\sum X_t^2 ,
\]
где использованы $\E\bh = \beta$ и $\D(\bh) = \sigma^2\big/\sum X_t^2$ из
\eqref{eq:noint}. Вычитая, получаем $\E\sum e_t^2 = (n-1)\sigma^2$, поэтому
несмещённая оценка дисперсии ошибок равна
\[
  s^2 = \frac{1}{n-1}\sum e_t^2 .
\]
Делитель $n-1$: оценивается один параметр.

\medskip
\textbf{б)} По \eqref{eq:noint} $\D(\bh) = \sigma^2\big/\sum X_t^2$, а её
оценка равна $s_{\bh}^2 = s^2\big/\sum X_t^2$.

\medskip
\textbf{в)} При предпосылке 3c оценка $\bh = \beta + \sum X_t\varepsilon_t\big/\sum X_t^2$
нормальна, поэтому
\[
  Z = \frac{\bh - \beta}{\sigma\big/\sqrt{\sum X_t^2}} \sim N(0,1).
\]
Кроме того, можно показать, что $V = (n-1)s^2/\sigma^2 \sim \chi^2(n-1)$ и не
зависит от $\bh$ (степеней свободы на единицу меньше, чем наблюдений, так как
оценён один параметр). Отношение
\[
  \frac{Z}{\sqrt{V/(n-1)}}
  = \frac{(\bh - \beta)\sqrt{\sum X_t^2}\big/\sigma}{s/\sigma}
  = \frac{\bh - \beta}{s_{\bh}}
\]
по определению распределения Стьюдента имеет распределение $t(n-1)$.

\medskip
\textbf{г)} Подходит набор данных из задачи~2.9, пункт~б): там
$\RSS/\TSS = 1{,}50 > 1$, а $1 - \ESS/\TSS = -3{,}26 < 0$. Все три патологии —
расхождение двух формул, значение больше единицы и значение меньше нуля —
возникают по одной причине: без свободного члена $\sum e_t \neq 0$, разложение
\eqref{eq:ss} нарушается на величину $2\overline{Y}\sum e_t$ согласно
\eqref{eq:decomp-noint}, и формулы \eqref{eq:r2} перестают быть эквивалентными.
\end{solution}

